% Chapter 3. Methods chapter: exact definitions match the final four-method Python study.
\chapter{FPLS, Krylov Methods, and Regularisation}
\label{chap:methods}

Chapter~\ref{chap:background} showed that functional linear regression leads
to a compact-operator equation whose direct inverse is unstable.  This
chapter develops the four regularised estimators used in the Monte Carlo
study: FPCR, raw FPLS, Arnoldi FPLS, and the CG--FPLS procedure
of \textcite{babii_carrasco_tsafack_2026}.  Particular attention is paid to a
terminological point that otherwise causes confusion.  Classical functional
PLS can be computed by a sequential score-and-deflation algorithm of NIPALS
type or represented as an alternative PLS (APLS) problem on a Krylov
subspace.  Raw and Arnoldi FPLS are two numerical implementations of this
same response-space APLS estimator.  Babii's CG--FPLS explores the same
nested Krylov spaces but uses a different finite-iteration objective.  The
methods are consequently closely related without being interchangeable in
every respect.

To align the theoretical comparison with the computation, the definitions
below use the same centring, discretisation, rank safeguards, and tuning
conventions.  Proofs of the
Krylov-space identities and variational characterisations are collected in
Appendix~\ref{app:mathematical-proofs}.

\section{The empirical operator problem}
\label{sec:methods-empirical-operator}

Let $X_i^c=X_i-\mu_X$ and $y_i=Y_i-\mu_Y$.  In an empirical application the
unknown means would normally be replaced by their sample counterparts.  The
controlled simulation designs impose $\mu_X=\mu_Y=0$, so the generated
curves and responses enter directly and no intercept is estimated.
Define the empirical sampling operator
$T_n:\mathcal H\rightarrow\mathbb R^n$ and its adjoint by
\begin{align}
    (T_nb)_i
    &=\langle X_i^c,b\rangle,
    \label{eq:methods-sampling-operator}\\
    T_n^*a
    &=\frac1n\sum_{i=1}^{n}a_iX_i^c,
    \qquad a\in\mathbb R^n,
    \label{eq:methods-adjoint-operator}
\end{align}
where $\langle a,c\rangle_n=n^{-1}a^\top c$ is used in response space.  The
empirical covariance operator and cross-covariance element are then
\begin{equation}
    \widehat K=T_n^*T_n,
    \qquad
    \widehat r=T_n^*y,
    \label{eq:methods-empirical-moments}
\end{equation}
so unregularised least squares would attempt to solve
$\widehat Kb=\widehat r$.  With population-centred curves, $\widehat K$ is
generated by $n$ rank-one operators and hence has rank at most $n$; when the
curves are centred by their sample mean, the sharper bound is $n-1$.  The
simulations use the known zero means, so the former bound applies.  In either
case, $\widehat K$ is finite rank and approximates a compact population
operator, so this equation cannot be used as a stable unrestricted inverse.
Unless $K$ is injective, all slope estimands below are understood as the
minimum-norm representatives in
$\mathcal H_{\mathrm{id}}=\operatorname{Null}(K)^\perp$; components in
$\operatorname{Null}(K)$ are not identified by the data.

The simulation evaluates functions on an equally spaced grid
$t_1,\ldots,t_T$ and uses the discrete inner product
\begin{equation}
    \langle f,g\rangle_T
    =\frac1T\sum_{\ell=1}^{T}f(t_\ell)g(t_\ell)
    =\frac1T f^\top g.
    \label{eq:methods-discrete-inner-product}
\end{equation}
If $X$ denotes the centred $n\times T$ matrix of curves, the operators in
\eqref{eq:methods-empirical-moments} are represented exactly as
\begin{equation}
    \widehat K=\frac{X^\top X}{nT},
    \qquad
    \widehat r=\frac{X^\top y}{n},
    \qquad
    T_nb=\frac{Xb}{T}.
    \label{eq:methods-grid-moments}
\end{equation}
These scaling factors matter: omitting $T^{-1}$ would change the numerical
operator and the stopping threshold.  All four estimators receive the same
sample, grid, zero-mean convention, and inner product.

Two residuals will be used repeatedly.  The response residual at $b$ is
$y-T_nb$, whereas the moment residual is
\begin{equation}
    e(b)=\widehat r-\widehat Kb
        =T_n^*(y-T_nb).
    \label{eq:methods-moment-residual}
\end{equation}
The latter is the response residual after it has been mapped back to function
space.  Although the two contain related information, minimising their norms
is not the same finite-dimensional optimisation problem.

\section{Functional principal component regression}
\label{sec:methods-fpcr}

FPCR is the spectral benchmark.  Let
$(\widehat\lambda_j,\widehat v_j)$ denote the empirical eigenpairs of
$\widehat K$, ordered by decreasing eigenvalue.  For a component count $m$,
the estimator is
\begin{equation}
    \widehat\beta_m^{\mathrm{FPCR}}
    =\sum_{j=1}^{m}
      \frac{\widehat v_j^\top\widehat r}
           {\widehat\lambda_j}\widehat v_j,
    \label{eq:methods-fpcr}
\end{equation}
where the coordinate eigenvectors are Euclidean-normalised.  Formula
\eqref{eq:methods-fpcr} is the least-squares solution restricted to the first
$m$ empirical principal directions.  The component count is therefore a
hard spectral cut-off: components beyond $m$ are assigned an inverse weight
of zero instead of the potentially large value
$1/\widehat\lambda_j$.
In the implementation, an empirical eigenvalue is retained only when it
exceeds
$\varepsilon_{\mathrm{mach}}T\max_j|\widehat\lambda_j|$; this is a numerical
rank safeguard, not an additional statistical tuning parameter.
The reported path is further capped at
$\min\{m_{\max},n-1,T-1\}$ before GCV, and at the number of numerically
positive eigenvalues.

The reported FPCR estimator chooses $m$ by response-space generalised
cross-validation (GCV).  If
$\widehat y_m=T_n\widehat\beta_m^{\mathrm{FPCR}}$, the criterion is
\begin{equation}
    \operatorname{GCV}_{Y}(m)
    =\frac{n^{-1}\|y-\widehat y_m\|_2^2}
           {(1-m/n)^2}.
    \label{eq:methods-response-gcv}
\end{equation}
The numerator rewards in-sample fit and the denominator penalises the
effective dimension.  This is response GCV, not a criterion based on the
moment residual.  Section~\ref{sec:methods-cg} introduces a separate
moment-space FPCR pilot used only to estimate the noise scale for CG stopping.

FPCR orders directions by predictor variance.  It is therefore expected to
be effective when the slope signal is concentrated in leading covariance
eigenfunctions.  It may be less parsimonious when a response-relevant
direction has modest predictor variance, which motivates supervised PLS
directions.

\section{Classical functional PLS and the NIPALS algorithm}
\label{sec:methods-nipals}

Partial least squares originated as a component method for multivariate
regression.  NIPALS, conventionally expanded as \emph{non-linear iterative
partial least squares}, is one of its classical score-and-deflation
algorithms \parencite{wold_et_al_1984,geladi_kowalski_1986}.  The present
model has one scalar response and therefore corresponds to PLS1.  Its
functional version replaces predictor vectors by functions and ordinary dot
products by the $L^2$ inner product.

To see the construction, start with residual curves
$X_i^{(0)}=X_i^c$ and residual responses $y_i^{(0)}=y_i$.  At component
$h=1,\ldots,m$, a NIPALS-type FPLS algorithm performs the following steps.
\begin{enumerate}
    \item Form and normalise the residual cross-covariance direction,
    \begin{equation}
        c_h=\frac1n\sum_{i=1}^{n}y_i^{(h-1)}X_i^{(h-1)},
        \qquad
        w_h=\frac{c_h}{\|c_h\|}.
        \label{eq:methods-nipals-weight}
    \end{equation}

    \item Project the residual predictor curves onto this direction to obtain
    the score
    \begin{equation}
        t_{ih}=\langle X_i^{(h-1)},w_h\rangle.
        \label{eq:methods-nipals-score}
    \end{equation}

    \item Regress both residual objects on the score.  The functional
    predictor loading and scalar response loading are
    \begin{equation}
        p_h=\frac{\sum_i t_{ih}X_i^{(h-1)}}{\sum_i t_{ih}^2},
        \qquad
        \ell_h=\frac{\sum_i t_{ih}y_i^{(h-1)}}{\sum_i t_{ih}^2}.
        \label{eq:methods-nipals-loadings}
    \end{equation}

    \item Deflate the predictor and response,
    \begin{equation}
        X_i^{(h)}=X_i^{(h-1)}-t_{ih}p_h,
        \qquad
        y_i^{(h)}=y_i^{(h-1)}-t_{ih}\ell_h.
        \label{eq:methods-nipals-deflation}
    \end{equation}
\end{enumerate}
At the first step, $c_1=\widehat r$, so the first PLS weight is the
normalised sample cross-covariance.  Subsequent weights are computed after
deflation and are therefore not simply displayed as raw powers of
$\widehat K$.  If $W_m a=\sum_{h=1}^{m}a_hw_h$ and
$P_m a=\sum_{h=1}^{m}a_hp_h$, the corresponding slope can be reconstructed
as
\begin{equation}
    \widehat\beta_m^{\mathrm{NIPALS}}
    =W_m(P_m^*W_m)^{-1}\ell,
    \qquad \ell=(\ell_1,\ldots,\ell_m)^\top,
    \label{eq:methods-nipals-slope}
\end{equation}
provided the small loading matrix is nonsingular.

The deflation formulas provide the classical algorithmic view: each score
explains the part of the response not captured by earlier scores.  Under the
usual full-rank conditions, NIPALS and SIMPLS yield the same fixed-$m$
response fit despite using different intermediate coordinates
\parencite{dejong_1993}.  The next section identifies that common fit directly
as an optimisation over a Krylov subspace.  NIPALS is therefore a reference
description, not an additional estimator in the simulation.

\section{APLS and the Krylov representation}
\label{sec:methods-apls}

\textcite{delaigle_hall_2012} established an alternative functional PLS
representation that exposes the operator geometry hidden by sequential
deflation.  Define the empirical Krylov subspace
\begin{equation}
    \widehat{\mathcal K}_m
    =\mathcal K_m(\widehat K,\widehat r)
    =\operatorname{span}\{\widehat r,\widehat K\widehat r,
      \ldots,\widehat K^{m-1}\widehat r\}.
    \label{eq:methods-empirical-krylov}
\end{equation}
Standard response-space FPLS is characterised by
\begin{equation}
    \widehat\beta_m^{\mathrm{APLS}}
    \in\arg\min_{b\in\widehat{\mathcal K}_m}
      \|y-T_nb\|_n^2.
    \label{eq:methods-apls-objective}
\end{equation}
Its first-order condition is
\begin{equation}
    \widehat r-\widehat K\widehat\beta_m^{\mathrm{APLS}}
    \perp\widehat{\mathcal K}_m.
    \label{eq:methods-apls-galerkin}
\end{equation}
The NIPALS scores span the image
$T_n\widehat{\mathcal K}_m$, so the sequential algorithm and
\eqref{eq:methods-apls-objective} produce the same fitted responses at a
common component count.  When the response risk is strictly convex on the
restricted space, their slope functions are equal as well.  This is the
precise similarity the present thesis uses: classical NIPALS-type FPLS and
APLS are algorithmic and operator-theoretic descriptions of the same
standard PLS1 fit.  A proof is given in
Proposition~\ref{prop:app-nipals-apls-equivalence}.

Every element of \eqref{eq:methods-empirical-krylov} is a polynomial in the
empirical covariance operator applied to the cross-covariance:
\begin{equation}
    b=p_{m-1}(\widehat K)\widehat r.
    \label{eq:methods-krylov-polynomial}
\end{equation}
Unlike FPCR, this space is response-informed because
$\widehat r=T_n^*y$ determines its starting direction.  The iteration does
not retain the $m$ largest-variance directions automatically.  It constructs
a polynomial approximation to the inverse in directions that are visible in
the empirical relationship between $X$ and $Y$.

Let $g_j=\widehat K^{j-1}\widehat r$ for $j=1,\ldots,m$ and write
$b=\sum_{j=1}^{m}a_jg_j$.  The normal equations for
\eqref{eq:methods-apls-objective} are
\begin{equation}
    H_m^{Y}a=h_m^{Y},
    \label{eq:methods-response-hankel-system}
\end{equation}
with entries
\begin{equation}
    (H_m^{Y})_{jk}
    =\langle\widehat r,
      \widehat K^{j+k-1}\widehat r\rangle,
    \qquad
    (h_m^{Y})_j
    =\langle\widehat r,
      \widehat K^{j-1}\widehat r\rangle.
    \label{eq:methods-response-hankel-entries}
\end{equation}
The matrix is Hankel because each entry depends only on $j+k$.  The same
system is obtained by regressing $y$ on the design columns
$T_ng_1,\ldots,T_ng_m$.  This representation reveals both the close
relationship with PLS and the numerical weakness of raw covariance powers.

\section{Raw FPLS}
\label{sec:methods-raw-fpls}

Raw FPLS implements \eqref{eq:methods-apls-objective} literally.  It forms
the matrix
\begin{equation}
    G_m=
    [\widehat r,\widehat K\widehat r,\ldots,
      \widehat K^{m-1}\widehat r]
    \label{eq:methods-raw-krylov-matrix}
\end{equation}
without rescaling or orthogonalising its columns, constructs
$Z_m=T_nG_m$, and solves
\begin{equation}
    (Z_m^\top Z_m)a=Z_m^\top y,
    \qquad
    \widehat\beta_m^{\mathrm{raw}}=G_ma.
    \label{eq:methods-raw-normal-equations}
\end{equation}
In exact arithmetic, this is simply APLS in the power basis.  It is included
because it is the most direct translation of the theoretical Krylov span and
because it exposes a numerical issue that can remain hidden when only final
prediction errors are reported.

Repeated multiplication by $\widehat K$ progressively emphasises the
leading empirical eigendirections.  To see this, expand
$\widehat r=\sum_j\widehat r_j\widehat v_j$.  Then
\begin{equation}
    \widehat K^k\widehat r
    =\sum_j\widehat\lambda_j^k
      \widehat r_j\widehat v_j.
    \label{eq:methods-power-eigen-expansion}
\end{equation}
As $k$ increases, ratios
$(\widehat\lambda_j/\widehat\lambda_1)^k$ contract whenever
$\widehat\lambda_j<\widehat\lambda_1$.  More precisely, let
$\widehat\lambda_\star$ be the largest empirical eigenvalue whose
eigenspace has a non-zero projection of $\widehat r$, and let
$P_\star$ denote that projection.  Then, in the finite empirical spectrum,
\[
 \frac{\widehat K^k\widehat r}{\|\widehat K^k\widehat r\|}
 \longrightarrow
 \frac{P_\star\widehat r}{\|P_\star\widehat r\|}
\]
whenever $\widehat\lambda_\star>0$.  Thus the normalised raw columns
eventually become nearly parallel.  The effect can be slow at modest $k$
when several active eigenvalues are close, but the condition number of
$G_m$ can nevertheless grow rapidly.  Moreover, whenever $Z_m$ has full
column rank,
$\kappa_2(Z_m^\top Z_m)=\kappa_2(Z_m)^2$ exactly; if it is rank deficient,
both sides are infinite under the usual extended convention.  Small
floating-point perturbations can then create large changes in $a$, despite
the fact that the intended exact subspace remains well defined.

The raw path is deliberately not repaired by a pseudoinverse or hidden
orthogonalisation in the simulation.  A component count is marked
invalid if its normal-equation solve fails or produces non-finite values.
Five-fold cross-validation is evaluated only on component counts valid in
the full sample and every fold.  This makes numerical failure part of the
empirical comparison rather than silently changing the estimator.

\section{Arnoldi FPLS and reorthogonalisation}
\label{sec:methods-arnoldi}

Arnoldi FPLS retains the response-space objective
\eqref{eq:methods-apls-objective} but changes the coordinates used for the
Krylov space.  Starting from
$q_1=\widehat r/\|\widehat r\|$, Arnoldi forms
$z=\widehat Kq_j$ and removes its projections onto the existing basis.  The
implementation uses two passes of classical Gram--Schmidt (CGS2):
\begin{align}
    z&\leftarrow z-Q_j(Q_j^\top z),
    \label{eq:methods-cgs-first-pass}\\
    z&\leftarrow z-Q_j(Q_j^\top z),
    \label{eq:methods-cgs-second-pass}\\
    q_{j+1}&=z/\|z\|,
    \label{eq:methods-arnoldi-normalise}
\end{align}
where $Q_j=[q_1,\ldots,q_j]$.  A second pass is inexpensive at the component
counts considered here and reduces the loss of orthogonality caused by
rounding.  The implementation terminates basis expansion when
\begin{equation}
    \|z\|_2\leq \varepsilon_{\mathrm{rank}}
                    \max\{\|\widehat K\|_\infty,\eta_{\mathrm{tiny}}\},
    \qquad \varepsilon_{\mathrm{rank}}=10^{-10},
    \label{eq:methods-arnoldi-rank-test}
\end{equation}
where $\eta_{\mathrm{tiny}}$ is the smallest positive normal floating-point
number.  Expansion also terminates when the candidate norm is non-finite.
The resulting columns span the same $\widehat{\mathcal K}_m$ as the raw
powers in exact arithmetic, as proved in
Proposition~\ref{prop:app-arnoldi-span}.

The displayed Arnoldi operations use the Euclidean norm of the grid
coordinates, exactly as in the code.  On an equally weighted grid this norm
differs from the discrete $L^2$ norm in
\eqref{eq:methods-discrete-inner-product} only by the common factor
$\sqrt T$.  The factor changes neither the generated span nor the
orthogonality and normalised-condition diagnostics.

For each available $m$, Arnoldi FPLS solves
\begin{equation}
    \widehat a_m
    \in\arg\min_{a\in\mathbb R^m}\|y-T_nQ_ma\|_n^2,
    \qquad
    \widehat\beta_m^{\mathrm{Arnoldi}}=Q_m\widehat a_m.
    \label{eq:methods-arnoldi-objective}
\end{equation}
The response-space least-squares problem is solved by QR decomposition,
not by forming another Gram matrix.  If
$Z_m^{A}=T_nQ_m=Q_{Z,m}R_{Z,m}$ denotes the reduced QR factorisation, the
response solve is accepted through component $j$ only while
\begin{equation}
    |(R_{Z,m})_{jj}|
    >\varepsilon_{\mathrm{rank}}
      \max\{|(R_{Z,m})_{11}|,\eta_{\mathrm{tiny}}\}.
    \label{eq:methods-arnoldi-response-rank-test}
\end{equation}
On failure of this response-rank test, or after a non-finite solve, the path
is frozen at its last valid fitted estimate for larger requested component
counts.  Thus, orthogonalisation protects the functional basis, QR protects
the subsequent response solve, and the two stated tests handle numerical
rank loss.  These are fixed computational safeguards rather than
data-driven regularisation choices.

Figure~\ref{fig:methods-krylov-geometry} illustrates the distinction for the
population geometry of Model~1.  At $m=12$, every pair among the normalised
raw powers has absolute inner product above $0.997$, and their condition
number is of order $10^{16}$.  The CGS2 Arnoldi basis has no off-diagonal
inner product exceeding the displayed threshold and an orthogonality defect
below $10^{-15}$.  The precise floating-point values can vary slightly across
linear-algebra platforms.  The figure does
not show different statistical information: both panels represent the same
Krylov space.  It shows why the choice of basis can determine whether that
space is numerically usable.

\begin{figure}[htbp]
    \centering
    \includegraphics[width=0.96\textwidth]{krylov_geometry.pdf}
    \caption[Raw and Arnoldi representations of the Krylov space]{Raw and CGS2-Arnoldi representations of the same
    $12$-dimensional population Krylov space in Model~1.  An edge marks an
    absolute pairwise inner product above $0.20$.  The raw power vectors are
    almost collinear, whereas reorthogonalisation yields an orthonormal
    coordinate system.  The illustration uses $J=100$ basis coefficients and
    a grid of $T=200$ points.}
    \label{fig:methods-krylov-geometry}
\end{figure}

At a fixed common $m$, a successful raw solve and a full-rank Arnoldi solve
must give the same exact-arithmetic response fit because both minimise the
same criterion over the same subspace.  Their selected estimates need not be
identical in the Monte Carlo study, however.  Each method is tuned by its own
five-fold cross-validation curve; finite-precision instability can alter the
raw curve, invalidate candidate dimensions, and lead to a different selected
component count.  The comparison therefore separates exact subspace
equivalence from the behaviour of complete, tuned numerical procedures.

\section{Conjugate-Gradient FPLS}
\label{sec:methods-cg}

The fourth estimator follows the construction of
\textcite{babii_carrasco_tsafack_2026}.  It searches over the same spaces
$\widehat{\mathcal K}_m$ but defines the finite-$m$ iterate by
\begin{equation}
    \widehat\beta_m^{\mathrm{CG}}
    \in\arg\min_{b\in\widehat{\mathcal K}_m}
      \|\widehat r-\widehat Kb\|^2.
    \label{eq:methods-cg-moment-objective}
\end{equation}
This is PLS measured in the norm of the empirical normal-equation residual.
Its first-order condition is
\begin{equation}
    \widehat r-\widehat K\widehat\beta_m^{\mathrm{CG}}
    \perp\widehat K\widehat{\mathcal K}_m,
    \label{eq:methods-cg-projection-condition}
\end{equation}
or, equivalently, $\widehat K e_m\perp\widehat{\mathcal K}_m$.  Compare this
with $e_m\perp\widehat{\mathcal K}_m$ in
\eqref{eq:methods-apls-galerkin}.  The subspace is the same, but the
projection condition is not.

The difference is also visible in raw coordinates.  With
$b=\sum_{j=1}^{m}a_jg_j$, the moment objective produces
\begin{equation}
    H_m^{M}a=h_m^{M},
    \label{eq:methods-moment-hankel-system}
\end{equation}
where
\begin{equation}
    (H_m^{M})_{jk}
    =\langle\widehat r,
      \widehat K^{j+k}\widehat r\rangle,
    \qquad
    (h_m^{M})_j
    =\langle\widehat r,
      \widehat K^{j}\widehat r\rangle.
    \label{eq:methods-moment-hankel-entries}
\end{equation}
Relative to \eqref{eq:methods-response-hankel-entries}, every operator power
has shifted by one.  Hence the NIPALS/APLS and CG estimates
generally differ before convergence.  Proposition
\ref{prop:app-response-moment-krylov} proves these two characterisations.

The estimator is computed without explicitly constructing either Hankel
matrix.  Initialise
\begin{equation}
    \widehat\beta_0=0,
    \qquad e_0=\widehat r,
    \qquad d_0=\widehat r.
\end{equation}
For $m=0,1,\ldots$, the released recurrence is
\begin{align}
    \alpha_m
    &=\frac{\langle e_m,\widehat K e_m\rangle}
            {\|\widehat Kd_m\|^2},
    \label{eq:methods-cg-alpha}\\
    \widehat\beta_{m+1}
    &=\widehat\beta_m+\alpha_md_m,
    \label{eq:methods-cg-beta-update}\\
    e_{m+1}
    &=e_m-\alpha_m\widehat Kd_m,
    \label{eq:methods-cg-residual-update}\\
    \gamma_m
    &=\frac{\langle e_{m+1},\widehat K e_{m+1}\rangle}
            {\langle e_m,\widehat K e_m\rangle},
    \label{eq:methods-cg-gamma}\\
    d_{m+1}
    &=e_{m+1}+\gamma_md_m.
    \label{eq:methods-cg-direction-update}
\end{align}
Its short recurrence avoids the explicit, ill-conditioned power-basis normal equations.  It should be noted that in
standard numerical-linear-algebra terminology, the residual-minimising
updates in \eqref{eq:methods-cg-alpha}--\eqref{eq:methods-cg-direction-update}
are the conjugate-residual recurrence, rather than the textbook energy-norm
CG recurrence used in Section~\ref{sec:methods-cg-versus-sd}.  The displayed
updates are well defined in exact arithmetic while their quadratic-form
numerator and squared-norm denominator are positive.  The implementation
symmetrises $\widehat K$ and freezes the remaining path at the last valid
iterate if either scalar is non-finite or non-positive, or if
\begin{equation}
    \|\widehat Kd_m\|_2
    \leq 10^{-12}
       \max\{\|\widehat K\|_\infty,\eta_{\mathrm{tiny}}\}
       \max\{\|d_m\|_2,\eta_{\mathrm{tiny}}\},
    \label{eq:methods-cg-breakdown-test}
\end{equation}
using the same $\eta_{\mathrm{tiny}}$ as above.  A non-finite updated iterate
or recurrence coefficient raises a floating-point error.  These are
numerical-breakdown guards, not the statistical stopping rule.

The theoretical discrepancy rule in
\textcite{babii_carrasco_tsafack_2026} uses the population noise scale
\begin{equation}
    \nu_n^\circ
    =\tau\left\{
      \frac{2\sigma^2\mathbb E\|X^c\|^2}{\delta n}
      \right\}^{1/2},
    \qquad
    m^\circ=\min\{m\geq0:\|e_m\|\leq\nu_n^\circ\}.
    \label{eq:methods-cg-ideal-threshold}
\end{equation}
Here $\sigma^2=\mathbb E\varepsilon^2$ in the homoskedastic simulation
design.  This ideal rule states the theoretical benchmark.  Because $\sigma^2$ and
$\mathbb E\|X^c\|^2$ are not supplied to an estimator, the implemented
workflow instead uses the plug-in scale
\begin{equation}
    \widehat\nu_n
    =\tau\left\{
      \frac{2\widehat\sigma^2\widehat{\mathbb E}\|X^c\|^2}
           {\delta n}\right\}^{1/2},
    \qquad
    \widehat{\mathbb E}\|X^c\|^2
    =\frac1n\sum_{i=1}^{n}\|X_i^c\|^2.
    \label{eq:methods-cg-threshold}
\end{equation}
and selects
\begin{equation}
    \widehat m
    =\min\left\{m\geq1:
      \|\widehat r-\widehat K\widehat\beta_m^{\mathrm{CG}}\|
      \leq\widehat\nu_n\right\},
    \label{eq:methods-cg-stopping-rule}
\end{equation}
with the maximum permitted component count used if the threshold is not
reached.  Here $\tau>1$ is a safety multiplier, $\delta\in(0,1)$ controls
the probability bound underlying the rule, and $\widehat\sigma^2$ is the
mean squared response residual from an internal FPCR pilot.  The pilot
dimension is chosen by moment-space GCV solely to estimate the noise level;
it does not enter the method comparison as a separate estimator.

\section{Why conjugate directions improve on steepest descent}
\label{sec:methods-cg-versus-sd}

The following comparison supplies geometric intuition for conjugate
directions; it does not assert equivalence with the released CG--FPLS
recurrence, which uses the conjugate-residual updates specified above.
Consider an ordinary symmetric positive-definite quadratic,
\begin{equation}
    f(x)=\frac12x^\top Ax-b^\top x,
    \qquad A=A^\top>0.
    \label{eq:methods-spd-quadratic}
\end{equation}
Steepest descent moves in the negative-gradient direction
$r_k=b-Ax_k$ and, even with exact line search, often zig-zags across a narrow
elliptical valley.  Its progress deteriorates as the condition number
$\kappa_2(A)$ increases.  Standard CG instead constructs directions satisfying
$p_j^\top Ap_k=0$ for $j\neq k$.  Minimisation along a new direction does
not undo minimisation along earlier conjugate directions, and exact-arithmetic
CG terminates in at most the dimension of the system
\parencite{hestenes_stiefel_1952}.

Figure~\ref{fig:methods-cg-sd-geometry} uses the same initial point, quadratic,
and convergence tolerance for both methods.  The matrix has condition number
$50$.  Steepest descent is given its optimal exact line-search step at every
iteration, yet takes $59$ iterations; CG reaches the two-dimensional minimiser
in two.  The contour and surface views depict the identical trajectories from
different perspectives.

\begin{figure}[htbp]
    \centering
    \begin{subfigure}[t]{0.49\textwidth}
        \centering
        \includegraphics[width=\linewidth]{quadratic_trajectory_2d.pdf}
        \caption{Contour view.}
    \end{subfigure}
    \hfill
    \begin{subfigure}[t]{0.49\textwidth}
        \centering
        \includegraphics[width=\linewidth]{quadratic_surface_3d.pdf}
        \caption{Quadratic-surface view.}
    \end{subfigure}
    \caption[Steepest descent and conjugate-gradient trajectories]{Steepest descent and conjugate gradient on the same
    two-dimensional positive-definite quadratic with $\kappa_2(A)=50$.
    Steepest descent uses exact line search.  The objective gap is measured
    relative to the common minimiser.}
    \label{fig:methods-cg-sd-geometry}
\end{figure}

In this limited role, the illustration motivates conjugate directions without
predicting the statistical ranking of the estimators.  That ranking also
depends on sampling error, ill-posedness, the restricted objective, and
tuning---the distinctions assembled after the stopping rule is considered.

\section{Early stopping as regularisation}
\label{sec:methods-early-stopping}

In exact finite-dimensional arithmetic, continued Krylov iteration can
eventually solve the estimable part of the empirical equation.  That is not
the statistical objective.  The right-hand side $\widehat r$ and operator
$\widehat K$ contain sampling noise, and late iterates increasingly fit
directions associated with small empirical eigenvalues.  The initial
iterates are typically smoother and more stable; the sequence may then enter
a semi-convergent regime in which the algebraic residual continues to fall
while slope error begins to rise.  Iteration number is therefore a
regularisation parameter rather than only a measure of computational effort
\parencite{hanke_1995,engl_hanke_neubauer_1996}.

The discrepancy principle formalises this idea by comparing the moment
residual with a sampling-noise scale.  If the residual is larger than that
scale, the iterate has not yet solved the reliable part of the empirical
equation.  Once it is comparable with sampling uncertainty, further
reduction is more likely to reproduce noise than stable signal.  Under the
source and moment conditions imposed by
\textcite{babii_carrasco_tsafack_2026}, the ideal rule in
\eqref{eq:methods-cg-ideal-threshold} adapts to unknown slope regularity and
yields convergence rates for both the unweighted slope norm and
covariance-weighted norms.  The source condition has the generic form
$\beta=K^\mu g$: a larger $\mu$ means that the slope is more concentrated in
well-identified covariance directions.  The practical rule
\eqref{eq:methods-cg-stopping-rule} replaces the unknown scale by a pilot
estimate.  Its behaviour is evaluated as part of the complete simulation
procedure; the theorem for the ideal threshold should not be read as an
automatic guarantee for this particular plug-in pilot without additional
consistency and probability-control arguments.

The discrepancy rule above is calibrated for estimation.  Inference instead
uses the stricter residual requirement derived in
Section~\ref{sec:background-estimation-inference}, and therefore follows the
same path to a later stopping index.  For a literal hypothesis
$H_0:\beta=b$, either $K$ must be injective or both $\beta$ and $b$ must be
restricted to $\mathcal H_{\mathrm{id}}$; otherwise the data identify only
the moment hypothesis $H_0:K\beta=Kb$.

\section{Equivalence boundaries and tuning design}
\label{sec:methods-comparison-framework}

The comparison is organised by three decisions: which subspace is searched,
which norm is minimised within that space, and how its dimension is chosen.
Table~\ref{tab:methods-equivalence} makes these dependencies explicit; a
common component count implies a common estimate only when both the subspace
and restricted objective coincide.

\begin{table}[!htbp]
    \centering
    \small
    \caption[Roles and equivalence boundaries of the methods]{Roles and exact equivalence boundaries of the methods discussed
    in this chapter.}
    \label{tab:methods-equivalence}
    \begin{tabular}{@{}
        >{\raggedright\arraybackslash}p{2.7cm}
        >{\raggedright\arraybackslash}p{3.3cm}
        >{\raggedright\arraybackslash}p{3.6cm}
        >{\raggedright\arraybackslash}p{4.3cm}@{}}
        \toprule
        Procedure & Search space & Finite-$m$ criterion & Relationship and tuning \\
        \midrule
        FPCR
        & Leading eigenspace of $\widehat K$
        & Response least squares
        & Spectral rather than response-informed; response GCV selects $m$ \\

        NIPALS-type FPLS
        & Sequential scores spanning $T_n\widehat{\mathcal K}_m$
        & Response least squares
        & Classical PLS1 reference; not an extra reported estimator \\

        Raw FPLS
        & Raw powers spanning $\widehat{\mathcal K}_m$
        & Response least squares
        & Same fixed-$m$ APLS fit as NIPALS under full rank; five-fold CV \\

        Arnoldi FPLS
        & CGS2 basis spanning $\widehat{\mathcal K}_m$
        & Response least squares solved by QR
        & Same exact fixed-$m$ fit as raw FPLS; five-fold CV applied separately \\

        CG--FPLS
        & Recursive directions spanning $\widehat{\mathcal K}_m$
        & Moment-residual least squares
        & Same nested spaces but generally a different finite-$m$ iterate;
          discrepancy stopping \\
        \bottomrule
    \end{tabular}
\end{table}

Tuning is part of each complete estimator and is therefore retained rather
than artificially equalised.  Separate tuning means that differences in
selected complexity are outcomes to interpret alongside prediction error,
slope error, and conditioning.  In particular, the raw--Arnoldi contrast asks
whether coordinate instability changes an otherwise common APLS target,
whereas the Arnoldi--CG contrast changes the finite-step objective itself.

These contrasts determine the Monte Carlo design in
Chapter~\ref{chap:simulation-design}, which holds the observations and
candidate path common while preserving each procedure's objective, tuning
rule, and numerical representation.
